%% notes-tex/025-additive-baselines/sections/4-sees.tex
%% Long version only: the reading of the 010 model code, carried over from
%% notes-tex/010-additive-baselines, and the comparability audit.

\section{What each model actually sees\secstatus{todo}}
\label{sec:sees}

The comparison is only meaningful if the models are given the same information.
Reading the 010 code settles that, and the answer is stronger than expected: every
model in the ladder, the transformer included, receives exactly one thing per
record, the set of three perturbed genes.

\subsection{The transformer's per-sample input\secstatus{todo}}

The dataset is built with the \texttt{Perturbation} graph processor. That processor
writes the perturbed gene names, their indices, a boolean mask and the label into
each sample, and nothing else (\texttt{graph\_processor.py}, lines 1984 to 2058).
There is no per-gene feature matrix: with an empty \texttt{node\_embeddings} list
in the config, the gene feature tensor of the cell graph has shape $6607 \times 0$
(\texttt{cell\_data.py}, line 53). Gene identity is carried by a free lookup table
of 6{,}607 rows and 180 columns
(\texttt{equivariant\_cell\_graph\_transformer.py}, line 1989), initialized randomly
and fit by gradient descent. No sequence, expression, ontology or other precomputed
representation enters. The model is given gene identities, the same as a one-hot
design matrix, and learns its own 180-dimensional vector for each.

\subsection{The encoder does not depend on the strain\secstatus{todo}}
\label{sec:encoder}

The forward pass builds its input as \texttt{gene\_embedding(torch.arange(N))},
prepends a CLS token, and unsqueezes to shape $[1, N{+}1, 180]$. That leading
dimension is $1$, not the batch size. The eight transformer layers then run on this
single tensor, which is assembled entirely from parameters and contains nothing
from \texttt{batch}. Writing $E$ for the embedding table and $T$ for the
eight-layer encoder,
%
\begin{equation}
  Z \;=\; T(E) \in \mathbb{R}^{N \times 180},
  \qquad
  h_{\mathrm{CLS}} \in \mathbb{R}^{180},
  \label{eq:encoder}
\end{equation}
%
and $Z$ and $h_{\mathrm{CLS}}$ are the same for every record in the dataset. The
repository states this itself, in the docstring of
\texttt{PerturbationGraphPropagation}: ``The encoder runs at batch 1 on the
WILDTYPE graph, so \texttt{H\_genes} and \texttt{h\_CLS} are identical for every
strain; the only strain-dependent step is
\texttt{EquivariantPerturbationTransform}.''

So the eight encoder layers, which hold 3{,}129{,}120 of the model's 4{,}774{,}861
parameters, perform no per-sample computation. They map one free parameter block to
another. Whatever role they play, it is not to process the strain. What remains
after Eq.~\ref{eq:encoder} is the strain-dependent part, Eq.~\ref{eq:cgt}, and that
is a set function on the three perturbed gene identities. The models in the ladder
differ in how expressive the set function is, not in what they are allowed to see.
Figure~\ref{fig:forward} draws the pass with the two entry points marked.

%% SOURCE: notes/experiments.025-solid-growth.mermaid.cgt-forward-pass.md, rendered
%% by notes/assets/publish/scripts/mermaid_pdf.sh and scaled by `make diagrams`;
%% the structure is read from torchcell/models/equivariant_cell_graph_transformer.py.
\tcfig{figures/cgt_forward_pass.pdf}{The forward pass of the 010 configuration, and
the additive null beside it. Yellow is a parameter block, orange a computation,
purple what a record contributes and what comes out, red what exists only during
training, blue the null on the same input. The record enters once, as three gene
indices, at the perturbation transform; the graphs enter once, through the layer-1
attention penalty, and reach the prediction only through the trained
weights.}{fig:forward}

\subsection{Where the nine graphs act, and where they do not\secstatus{todo}}
\label{sec:graphs}

The nine adjacency matrices enter the model in exactly one place. During training,
a Kullback-Leibler penalty pulls one attention head of layer 1 toward each
row-normalized adjacency $\tilde{A}^{(k)}$,
%
\begin{equation}
  \mathcal{L}_{\mathrm{graph}} \;=\; \sum_{k=1}^{9} c_k
  \sum_{i=1}^{N} \sum_{j=1}^{N}
  \tilde{A}^{(k)}_{ij} \log \frac{\tilde{A}^{(k)}_{ij}}{\alpha^{(k)}_{ij}},
  \qquad
  \alpha^{(k)} = \mathrm{softmax}\!\left(\frac{QK^{\top}}{\sqrt{d}}\right)_{\!h_k},
  \label{eq:graphreg}
\end{equation}
%
and that is the only operation that reads them.

\notesec{Not in the prediction}
Three mechanisms could put adjacency into the computation that yields $\hat{y}$: a
head mask on attention, the perturbation-propagation block, and an adjacency bias
inside the attention logits. All three are off in the 010 configuration. Given a
trained checkpoint, no adjacency tensor is an input to any operation between the
gene indices and the prediction, and Section~\ref{sec:verify} checks this directly:
setting the penalty weight to zero on a trained checkpoint moves the predictions by
at most $1.3 \times 10^{-6}$.

\notesec{In the weights}
The penalty is computed inside the same forward call as the prediction, from the
same pre-dropout attention tensor, and nothing detaches it. Its gradient reaches the
layer-1 query and key projections, through them every layer-0 parameter, and
through those the gene embedding table. So $Z = T(E)$ of Eq.~\ref{eq:encoder}, the
per-gene vector the prediction is read from, is shaped by the nine graphs during
training even though its value at inference comes off fixed weights. The graphs act
on the parameters, not on the function from checkpoint to prediction. An ablation
that zeroes the coefficient on a trained checkpoint tests only the second.

\notesec{How large the penalty was}
The per-graph weight $c_k$ in Eq.~\ref{eq:graphreg} is a product of three factors:
the per-head $\lambda_k = 10^{-3}$, a global scale of $1.0$ read from
\texttt{regression\_task.loss.graph\_regularization.lambda}, and an edge
normalization meant to divide by the mean degree. That last factor is computed as
\texttt{len(A.nonzero()[0])} summed over the nine matrices. On a dense
\texttt{torch} tensor \texttt{A.nonzero()} has shape $[\mathrm{nnz}, 2]$, so
\texttt{[0]} is one index pair of length two, and the sum over nine graphs is $18$
rather than the 2.5 million edges intended. Dividing by $18/6607$ multiplies by
$367.06$,
%
\begin{equation}
  c_k \;=\; 10^{-3} \times 1.0 \times \frac{6607}{18} \;=\; 0.367 ,
  \label{eq:effcoef}
\end{equation}
%
applied to a divergence summed over all $6{,}607$ rows. The 010 runs' logs record
the consequence: against a point loss of $0.81$ the graph term is $5{,}740$, and the
trainer's diagnostic \texttt{norm\_weighted\_graph\_reg}, the graph term divided by
the total loss, reads $0.99986$ for all three checkpoints. The optimizer spent 99.99
percent of its objective on the graph penalty. Figure~\ref{fig:penalty} shows the
three loss terms and that share over the whole of training, and
Section~\ref{sec:graphablation} reports what the penalty bought.

The dominance was not what produced the accuracy. The 025 replication of this
configuration, GH 1598, ran after the edge normalization had been corrected, so its
graph term is $0.34$ against a point loss of $0.80$, a share of 0.29 to 0.33 rather
than 0.9999, and it reached validation Pearson 0.446 at epoch 14 against 0.447 to
0.462 for the three 010 checkpoints.

%% SOURCE: experiments/010-kuzmin-tmi/scripts/graph_penalty_vs_loss.py ->
%% results/graph_penalty_vs_loss_history.csv, pulled from the four runs' W&B
%% histories; summary in results/graph_penalty_vs_loss_summary.json
\tcfig{figures/graph_penalty_vs_loss.pdf}{The training objective of the three 010
runs and of the 025 replication, per epoch. \textbf{a}, the point term, the
Sinkhorn term at its weight of 0.1, and the weighted graph penalty for the 010 runs,
on a log axis: the penalty sits four orders of magnitude above the point loss from
epoch 5 on. \textbf{b}, the same three terms for GH 1598, where the corrected edge
normalization puts the penalty below the point loss. \textbf{c}, the graph term's
share of the total loss, recomputed from the logged terms: 0.9999 throughout for the
010 runs, 0.32 for the replication. \textbf{d}, validation Pearson per epoch with
the best-Pearson epoch marked; the replication reaches the 010 band with the penalty
at a third of the objective.}{fig:penalty}

All nine heads carried the penalty. The configuration names the first two graphs
\texttt{physical} and \texttt{regulatory}, which is what the cell graph called them
at the December 2025 commit these runs used; the suffixed names date from
2026-06-27. The runs log the edge-recovery keys \texttt{physical\_L1\_H0} and
\texttt{regulatory\_L1\_H1}, which the trainer emits only for a graph it matched.

\subsection{Is the comparison fair\secstatus{todo}}
\label{sec:comparability}

Six things have to match before a baseline number can sit in the same table as a
transformer number. Table~\ref{tab:comparability} names them and gives the verdict
on each.

%% SOURCE: each row is established in the text it names; the normalization
%% constants are read from the runs' logged transform config and from
%% experiments/025-solid-growth/results/label_normalization_constants.json.
\begin{table}[htbp]
\centering
\footnotesize
\caption{What has to match for the comparison to be fair. Row 6 is the one
mismatch, and it holds for arm R only. It is a hygiene defect in the 010
configuration rather than a source of advantage, and arm Q fits its constants on
the training part.}
\label{tab:comparability}
\begin{tabular}{llp{0.62\textwidth}}
\toprule
\# & Verdict & What has to match \\
\midrule
1 & matches & \textbf{Records and split.} All models use the same 376{,}732 records
    and, per arm, the same partition (Table~\ref{tab:arms}). \\
2 & matches & \textbf{Label.} All predict \texttt{gene\_interaction}, the published
    adjusted $\tau$ of Eq.~\ref{eq:tau}. \\
3 & matches & \textbf{Reported split.} Test is reported for every baseline row, and
    validation is what every row selected on. The transformer rows that are
    validation maxima say so in their label. \\
4 & matches & \textbf{Metric definition and denominator.} Pearson, Spearman and
    mean squared error over the full held-out part in raw $\tau$ units. \\
5 & matches & \textbf{Information available.} The identity of the three perturbed
    genes, and nothing else, for every model (Sections~\ref{sec:encoder}
    and~\ref{sec:cgt}). \\
6 & arm R only & \textbf{Label normalization.} The arm R transformer trains on
    z-scored labels whose mean and standard deviation were computed from all
    376{,}732 records, so two held-out scalars touch a training-time constant; the
    baselines use the training mean only. Pearson is invariant to the two
    constants, so this cannot move the reported metric. Arm Q fits on its 301{,}236
    training records. \\
\bottomrule
\end{tabular}
\end{table}

\notesec{Label normalization: what it can and cannot move}
The transformer is trained on the z-scored label $(y - \mu) / \sigma$. In the 010
configuration and arm R the two constants are the all-record statistics,
$\mu = -0.008024324$ and $\sigma = 0.063263549$, where the training split alone
gives $-0.007688739$ and $0.063186255$. Pearson and Spearman are invariant to an
affine map of the labels, so no choice of $\mu$ and $\sigma$ changes the correlation
a given model would reach. What the constants set is the scale of the point loss
relative to the other terms of Eq.~\ref{eq:loss}: a different $\sigma$ multiplies
the standardized mean squared error by $(\sigma_{\mathrm{all}} /
\sigma_{\mathrm{train}})^2 = 1.0024$, and a different $\mu$ is a constant offset the
readout bias absorbs. Arm Q names its fit population, \texttt{transforms.fit\_on:
train}, and its constants are $\mu = -0.007843739$ and $\sigma = 0.063776418$; both
sets are written by
\file{experiments/025-solid-growth/scripts/label_normalization_constants.py}.

\notesec{The loss functions differ, deliberately}
The baselines minimize squared error. The transformer minimizes
%
\begin{equation}
  \mathcal{L} \;=\; \lambda_p \mathcal{L}_{\mathrm{MSE}}
  \;+\; \lambda_d\, S_{\epsilon}\!\left(
    \tfrac{1}{B}\textstyle\sum_b \delta_{\hat{y}_b},\;
    \tfrac{1}{B}\textstyle\sum_b \delta_{y_b}\right)
  \;+\; \lambda_g \mathcal{L}_{\mathrm{graph}},
  \label{eq:loss}
\end{equation}
%
with $\lambda_p = 1$, $\lambda_d = 0.1$, $\lambda_g = 1$. The middle term is a
debiased Sinkhorn divergence with $p = 2$ and blur $0.05$ between the empirical
measure of the batch's predictions and that of the batch's targets, computed
unpaired, which shapes the marginal distribution of the predictions rather than
their per-record accuracy. The third term is the graph penalty of
Eq.~\ref{eq:graphreg}. This is a difference in training objective, not in the
evaluation, and it is part of what the transformer is. A distributional term can
move Pearson and mean squared error in opposite directions, so the two metrics are
read together rather than either alone.

\notesec{Model size}

%% SOURCE: parameter counts read from
%% torchcell/models/equivariant_cell_graph_transformer.py at the lines given
\begin{table}[htbp]
\centering
\footnotesize
\caption{Parameter budget. Two thirds of the transformer sits in the encoder, which
does no per-sample work, and a quarter in the embedding table. The
strain-dependent computation is 456{,}301 parameters, under a tenth of the model.}
\label{tab:params}
\begin{tabular}{lrl}
\toprule
Component & Parameters & Role \\
\midrule
Gene embedding table & 1{,}189{,}260 & the only place gene identity lives \\
Encoder, 8 layers & 3{,}129{,}120 & strain-independent, Section~\ref{sec:encoder} \\
Perturbation transform & 391{,}140 & the strain-dependent step \\
Readout head & 65{,}161 & pooling and MLP \\
CLS token & 180 & constant \\
\midrule
Total & 4{,}774{,}861 & \\
\midrule
B1 additive ridge & 4{,}353 & one scalar per gene, plus intercept \\
B5 embedding MLP & 623{,}745 & 128 per gene, plus a two-layer network \\
\bottomrule
\end{tabular}
\end{table}
